\section{Projective space}
\label{section-projective-space}

\noindent
Projective space is one of the fundamental objects studied in
algebraic geometry. In this section we just give its construction
as $\text{Proj}$ of a polynomial ring. Later we will discover many
of its beautiful properties.

\begin{lemma}
\label{lemma-projective-space}
Let $S = \mathbf{Z}[T_0, \ldots, T_n]$ with $\deg(T_i) = 1$.
The scheme
$$
\mathbf{P}^n_{\mathbf{Z}} = \text{Proj}(S)
$$
represents the functor which associates to a scheme $Y$ the pairs
$(\mathcal{L}, (s_0, \ldots, s_n))$ where
\begin{enumerate}
\item $\mathcal{L}$ is an invertible $\mathcal{O}_Y$-module, and
\item $s_0, \ldots, s_n$ are global sections of $\mathcal{L}$
which generate $\mathcal{L}$
\end{enumerate}
up to the following equivalence:
$(\mathcal{L}, (s_0, \ldots, s_n)) \sim
(\mathcal{N}, (t_0, \ldots, t_n))$ $\Leftrightarrow$ there exists
an isomorphism $\beta : \mathcal{L} \to \mathcal{N}$
with $\beta(s_i) = t_i$ for $i = 0, \ldots, n$.
\end{lemma}

\begin{proof}
This is a special case of Lemma \ref{lemma-apply} above.
Namely, for any graded ring $A$ we have
\begin{eqnarray*}
\Mor_{graded rings}(\mathbf{Z}[T_0, \ldots, T_n], A)
& = &
A_1 \times \ldots \times A_1 \\
\psi & \mapsto & (\psi(T_0), \ldots, \psi(T_n))
\end{eqnarray*}
and the degree $1$ part of $\Gamma_*(Y, \mathcal{L})$ is
just $\Gamma(Y, \mathcal{L})$.
\end{proof}

\begin{definition}
\label{definition-projective-space}
The scheme
$\mathbf{P}^n_{\mathbf{Z}} = \text{Proj}(\mathbf{Z}[T_0, \ldots, T_n])$
is called {\it projective $n$-space over $\mathbf{Z}$}.
Its base change $\mathbf{P}^n_S$ to a scheme $S$ is called
{\it projective $n$-space over $S$}. If $R$ is a ring the base change
to $\Spec(R)$ is denoted $\mathbf{P}^n_R$ and called
{\it projective $n$-space over $R$}.
\end{definition}

\noindent
Given a scheme $Y$ over $S$
and a pair $(\mathcal{L}, (s_0, \ldots, s_n))$ as in
Lemma \ref{lemma-projective-space}
the induced morphism to $\mathbf{P}^n_S$ is denoted
$$
\varphi_{(\mathcal{L}, (s_0, \ldots, s_n))} :
Y \longrightarrow \mathbf{P}^n_S
$$
This makes sense since the pair defines a morphism into
$\mathbf{P}^n_{\mathbf{Z}}$ and we already have the structure
morphism into $S$ so combined we get a morphism into
$\mathbf{P}^n_S = \mathbf{P}^n_{\mathbf{Z}} \times S$.
Via Lemma \ref{lemma-relative-proj-base-change}
we can identify $\mathbf{P}^n_S$ with a relative Proj of
the graded $\mathcal{O}_S$-algebra $\mathcal{O}_S[T_0, \ldots, T_n]$
and we can characterize $\varphi_{(\mathcal{L}, (s_0, \ldots, s_n))}$
as the unique $S$-morphism such that
$$
\mathcal{L} =
\varphi_{(\mathcal{L}, (s_0, \ldots, s_n))}^*\mathcal{O}_{\mathbf{P}^n_S}(1)
\quad
\text{and}
\quad
s_i = \varphi_{(\mathcal{L}, (s_0, \ldots, s_n))}^*T_i
$$
where we think of $T_i$ as a global section of
$\mathcal{O}_{\mathbf{P}^n_S}(1)$ via the map $\psi$ of
Lemma \ref{lemma-glue-relative-proj-twists}; details omitted.

\begin{lemma}
\label{lemma-standard-covering-projective-space}
Projective $n$-space over $\mathbf{Z}$ is covered by
$n + 1$ standard opens
$$
\mathbf{P}^n_{\mathbf{Z}} =
\bigcup\nolimits_{i = 0, \ldots, n} D_{+}(T_i)
$$
where each $D_{+}(T_i)$ is isomorphic to $\mathbf{A}^n_{\mathbf{Z}}$
affine $n$-space over $\mathbf{Z}$.
\end{lemma}

\begin{proof}
This is true because
$\mathbf{Z}[T_0, \ldots, T_n]_{+} = (T_0, \ldots, T_n)$ and
since
$$
\Spec
\left(
\mathbf{Z}
\left[\frac{T_0}{T_i}, \ldots, \frac{T_n}{T_i}
\right]
\right)
\cong
\mathbf{A}^n_{\mathbf{Z}}
$$
in an obvious way.
\end{proof}

\begin{lemma}
\label{lemma-projective-space-separated}
Let $S$ be a scheme.
The structure morphism $\mathbf{P}^n_S \to S$ is
\begin{enumerate}
\item separated,
\item quasi-compact,
\item satisfies the existence and uniqueness parts of the valuative criterion,
and
\item universally closed.
\end{enumerate}
\end{lemma}

\begin{proof}
All these properties are stable under base change (this is clear for the
last two and for the other two see
Schemes, Lemmas
\ref{schemes-lemma-separated-permanence} and
\ref{schemes-lemma-quasi-compact-preserved-base-change}).
Hence it suffices to prove them for the morphism
$\mathbf{P}^n_{\mathbf{Z}} \to \Spec(\mathbf{Z})$.
Separatedness is Lemma \ref{lemma-proj-separated}. Quasi-compactness follows
from Lemma \ref{lemma-standard-covering-projective-space}.
Existence and uniqueness of the valuative criterion follow from
Lemma \ref{lemma-proj-valuative-criterion}.
Universally closed follows from the above and
Schemes, Proposition \ref{schemes-proposition-characterize-universally-closed}.
\end{proof}

\begin{remark}
\label{remark-missing-finite-type}
What's missing in the list of properties above? Well to be sure the property
of being of finite type. The reason we do not list this here is that we have
not yet defined the notion of finite type at this point. (Another property
which is missing is ``smoothness''. And I'm sure there are many more you can
think of.)
\end{remark}

\begin{lemma}[Segre embedding]
\label{lemma-segre-embedding}
Let $S$ be a scheme. There exists a closed immersion
$$
\mathbf{P}^n_S \times_S \mathbf{P}^m_S
\longrightarrow
\mathbf{P}^{nm + n + m}_S
$$
called the {\it Segre embedding}.
\end{lemma}

\begin{proof}
It suffices to prove this when $S = \Spec(\mathbf{Z})$.
Hence we will drop the index $S$ and work in the absolute setting.
Write $\mathbf{P}^n = \text{Proj}(\mathbf{Z}[X_0, \ldots, X_n])$,
$\mathbf{P}^m = \text{Proj}(\mathbf{Z}[Y_0, \ldots, Y_m])$,
and
$\mathbf{P}^{nm + n + m} =
\text{Proj}(\mathbf{Z}[Z_0, \ldots, Z_{nm + n + m}])$.
In order to map into $\mathbf{P}^{nm + n + m}$ we have to
write down an invertible sheaf $\mathcal{L}$ on the left hand
side and $(n + 1)(m + 1)$ sections $s_i$ which generate it.
See Lemma \ref{lemma-projective-space}.
The invertible sheaf we take is
$$
\mathcal{L} =
\text{pr}_1^*\mathcal{O}_{\mathbf{P}^n}(1)
\otimes
\text{pr}_2^*\mathcal{O}_{\mathbf{P}^m}(1)
$$
The sections we take are
$$
s_0 = X_0Y_0, \ s_1 = X_1Y_0, \ldots, \ s_n = X_nY_0,
\ s_{n + 1} = X_0Y_1, \ldots, \ s_{nm + n + m} = X_nY_m.
$$
These generate $\mathcal{L}$ since the sections $X_i$ generate
$\mathcal{O}_{\mathbf{P}^n}(1)$ and the sections $Y_j$ generate
$\mathcal{O}_{\mathbf{P}^m}(1)$. The induced morphism
$\varphi$ has the property that
$$
\varphi^{-1}(D_{+}(Z_{i + (n + 1)j})) = D_{+}(X_i) \times D_{+}(Y_j).
$$
Hence it is an affine morphism. The corresponding ring map in case
$(i, j) = (0, 0)$ is the map
$$
\mathbf{Z}[Z_1/Z_0, \ldots, Z_{nm + n + m}/Z_0]
\longrightarrow
\mathbf{Z}[X_1/X_0, \ldots, X_n/X_0, Y_1/Y_0, \ldots, Y_n/Y_0]
$$
which maps $Z_i/Z_0$ to the element $X_i/X_0$ for $i \leq n$ and
the element $Z_{(n + 1)j}/Z_0$ to the element $Y_j/Y_0$. Hence it
is surjective. A similar argument works for the other affine
open subsets. Hence the morphism $\varphi$ is a closed immersion
(see Schemes, Lemma \ref{schemes-lemma-closed-local-target} and
Example \ref{schemes-example-closed-immersion-affines}.)
\end{proof}

\noindent
The following two lemmas are special cases of more general results later, but
perhaps it makes sense to prove these directly here now. The first one is
a special case of Divisors, Lemma \ref{divisors-lemma-closed-subscheme-proj}.

\begin{lemma}
\label{lemma-closed-in-projective-space}
Let $R$ be a ring. Let $Z \subset \mathbf{P}^n_R$ be a closed subscheme.
Let
$$
I_d = \Ker\left(
R[T_0, \ldots, T_n]_d
\longrightarrow
\Gamma(Z, \mathcal{O}_{\mathbf{P}^n_R}(d)|_Z)\right)
$$
Then $I = \bigoplus I_d \subset R[T_0, \ldots, T_n]$ is
a graded ideal and $Z = \text{Proj}(R[T_0, \ldots, T_n]/I)$.
\end{lemma}

\begin{proof}
It is clear that $I$ is a graded ideal.
Set $Z' = \text{Proj}(R[T_0, \ldots, T_n]/I)$.
By Lemma \ref{lemma-surjective-graded-rings-generated-degree-1-map-proj}
we see that $Z'$ is a closed subscheme of $\mathbf{P}^n_R$.
To see the equality $Z = Z'$
it suffices to check on an standard affine open
$D_{+}(T_i)$. By renumbering the homogeneous coordinates we
may assume $i = 0$. Say $Z \cap D_{+}(T_0)$, resp.\ $Z' \cap D_{+}(T_0)$
is cut out by the ideal $J$, resp.\ $J'$ of $R[T_1/T_0, \ldots, T_n/T_0]$.
Then $J'$ is the ideal generated by the elements $F/T_0^{\deg(F)}$ where
$F \in I$ is homogeneous.
Suppose the degree of $F \in I$ is $d$. Since $F$ vanishes as a section
of $\mathcal{O}_{\mathbf{P}^n_R}(d)$ restricted to $Z$ we see that
$F/T_0^d$ is an element of $J$. Thus $J' \subset J$.

\medskip\noindent
Conversely, suppose that $f \in J$. If $f$ has total degree
$d$ in $T_1/T_0, \ldots, T_n/T_0$, then we can write
$f = F/T_0^d$ for some $F \in R[T_0, \ldots, T_n]_d$.
Pick $i \in \{1, \ldots, n\}$. Then $Z \cap D_{+}(T_i)$ is
cut out by some ideal $J_i \subset R[T_0/T_i, \ldots, T_n/T_i]$.
Moreover,
$$
J \cdot
R\left[
\frac{T_1}{T_0}, \ldots, \frac{T_n}{T_0},
\frac{T_0}{T_i}, \ldots, \frac{T_n}{T_i}
\right]
=
J_i \cdot
R\left[
\frac{T_1}{T_0}, \ldots, \frac{T_n}{T_0},
\frac{T_0}{T_i}, \ldots, \frac{T_n}{T_i}
\right]
$$
The left hand side is the localization of $J$ with respect to the
element $T_i/T_0$ and the right hand side is the localization of $J_i$
with respect to the element $T_0/T_i$. It follows that
$T_0^{d_i}F/T_i^{d + d_i}$ is an element of $J_i$ for some $d_i$
sufficiently large. This proves that $T_0^{\max(d_i)}F$ is an
element of $I$, because its restriction to each standard affine
open $D_{+}(T_i)$ vanishes on the closed subscheme
$Z \cap D_{+}(T_i)$. Hence $f \in J'$ and we conclude $J \subset J'$
as desired.
\end{proof}

\noindent
The following lemma is a special case of the more general
Properties, Lemmas \ref{properties-lemma-ample-quasi-coherent} or
\ref{properties-lemma-proj-quasi-coherent}.

\begin{lemma}
\label{lemma-quasi-coherent-projective-space}
Let $R$ be a ring.
Let $\mathcal{F}$ be a quasi-coherent sheaf on $\mathbf{P}^n_R$.
For $d \geq 0$ set
$$
M_d
=
\Gamma(\mathbf{P}^n_R,
\mathcal{F} \otimes_{\mathcal{O}_{\mathbf{P}^n_R}}
\mathcal{O}_{\mathbf{P}^n_R}(d))
=
\Gamma(\mathbf{P}^n_R, \mathcal{F}(d))
$$
Then $M = \bigoplus_{d \geq 0} M_d$ is a graded $R[T_0, \ldots, R_n]$-module
and there is a canonical isomorphism $\mathcal{F} = \widetilde{M}$.
\end{lemma}

\begin{proof}
The multiplication maps
$$
R[T_0, \ldots, R_n]_e \times M_d \longrightarrow M_{d + e}
$$
come from the natural isomorphisms
$$
\mathcal{O}_{\mathbf{P}^n_R}(e)
\otimes_{\mathcal{O}_{\mathbf{P}^n_R}}
\mathcal{F}(d)
\longrightarrow
\mathcal{F}(e + d)
$$
see Equation (\ref{equation-global-sections-module}). Let us construct the
map $c : \widetilde{M} \to \mathcal{F}$. On each of the standard affines
$U_i = D_{+}(T_i)$ we see that $\Gamma(U_i, \widetilde{M}) = (M[1/T_i])_0$
where the subscript ${}_0$ means degree $0$ part. An element of this
can be written as $m/T_i^d$ with $m \in M_d$. Since $T_i$ is a generator
of $\mathcal{O}(1)$ over $U_i$ we can always write
$m|_{U_i} = m_i \otimes T_i^d$ where $m_i \in \Gamma(U_i, \mathcal{F})$
is a unique section. Thus a natural guess is $c(m/T_i^d) = m_i$.
A small argument, which is omitted here, shows that this gives a
well defined map $c : \widetilde{M} \to \mathcal{F}$ if we can
show that
$$
(T_i/T_j)^d m_i|_{U_i \cap U_j} = m_j|_{U_i \cap U_j}
$$
in $M[1/T_iT_j]$.
But this is clear since on the overlap the generators $T_i$ and
$T_j$ of $\mathcal{O}(1)$ differ by the invertible function $T_i/T_j$.

\medskip\noindent
Injectivity of $c$. We may check for injectivity over the affine opens
$U_i$. Let $i \in \{0, \ldots, n\}$
and let $s$ be an element $s = m/T_i^d \in \Gamma(U_i, \widetilde{M})$
such that $c(m/T_i^d) = 0$. By the description of $c$ above this means that
$m_i = 0$, hence $m|_{U_i} = 0$. Hence $T_i^em = 0$ in $M$ for some
$e$. Hence $s = m/T_i^d = T_i^e/T_i^{e + d} = 0$ as desired.

\medskip\noindent
Surjectivity of $c$. We may check for surjectivity over the affine opens
$U_i$. By renumbering it suffices to check it over $U_0$.
Let $s \in \mathcal{F}(U_0)$.
Let us write $\mathcal{F}|_{U_i} = \widetilde{N_i}$ for some
$R[T_0/T_i, \ldots, T_0/T_i]$-module $N_i$, which is possible because
$\mathcal{F}$ is quasi-coherent. So $s$ corresponds to an element
$x \in N_0$. Then we have that
$$
(N_i)_{T_j/T_i} \cong (N_j)_{T_i/T_j}
$$
(where the subscripts mean ``principal localization at'')
as modules over the ring
$$
R\left[
\frac{T_0}{T_i}, \ldots, \frac{T_n}{T_i},
\frac{T_0}{T_j}, \ldots, \frac{T_n}{T_j}
\right].
$$
This means that for some large integer $d$ there exist elements
$s_i \in N_i$, $i = 1, \ldots, n$ such that
$$
s = (T_i/T_0)^d s_i
$$
on $U_0 \cap U_i$. Next, we look at the difference
$$
t_{ij} = s_i - (T_j/T_i)^d s_j
$$
on $U_i \cap U_j$, $0 < i < j$. By our choice of $s_i$ we know that
$t_{ij}|_{U_0 \cap U_i \cap U_j} = 0$. Hence there exists a large integer
$e$ such that $(T_0/T_i)^et_{ij} = 0$. Set $s_i' = (T_0/T_i)^es_i$,
and $s_0' = s$. Then we will have
$$
s_a' = (T_b/T_a)^{e + d} s_b'
$$
on $U_a \cap U_b$ for all $a, b$. This is exactly the condition that the
elements $s'_a$ glue to a global section
$m \in \Gamma(\mathbf{P}^n_R, \mathcal{F}(e + d))$.
And moreover $c(m/T_0^{e + d}) = s$ by construction. Hence $c$ is
surjective and we win.
\end{proof}

\begin{lemma}
\label{lemma-globally-generated-omega-twist-1}
Let $X$ be a scheme. Let $\mathcal{L}$ be an invertible sheaf
and let $s_0, \ldots, s_n$ be global sections of $\mathcal{L}$
which generate it. Let $\mathcal{F}$ be the kernel of the induced
map $\mathcal{O}_X^{\oplus n + 1} \to \mathcal{L}$.
Then $\mathcal{F} \otimes \mathcal{L}$ is globally generated.
\end{lemma}

\begin{proof}
In fact the result is true if $X$ is any locally ringed space.
The sheaf $\mathcal{F}$ is a finite locally free $\mathcal{O}_X$-module
of rank $n$. The elements
$$
s_{ij} = (0, \ldots, 0, s_j, 0, \ldots, 0, -s_i, 0, \ldots, 0)
\in \Gamma(X, \mathcal{L}^{\oplus n + 1})
$$
with $s_j$ in the $i$th spot and $-s_i$ in the $j$th spot map to zero
in $\mathcal{L}^{\otimes 2}$. Hence
$s_{ij} \in \Gamma(X, \mathcal{F} \otimes_{\mathcal{O}_X} \mathcal{L})$.
A local computation shows that these sections generate
$\mathcal{F} \otimes \mathcal{L}$.

\medskip\noindent
Alternative proof. Consider the morphism
$\varphi : X \to \mathbf{P}^n_\mathbf{Z}$ associated to
the pair $(\mathcal{L}, (s_0, \ldots, s_n))$. Since the pullback
of $\mathcal{O}(1)$ is $\mathcal{L}$ and since the pullback
of $T_i$ is $s_i$, it suffices to prove the lemma in the
case of $\mathbf{P}^n_\mathbf{Z}$. In this case the sheaf
$\mathcal{F}$ corresponds to the graded $S = \mathbf{Z}[T_0, \ldots, T_n]$
module $M$ which fits into the short exact sequence
$$
0 \to M \to S^{\oplus n + 1} \to S(1) \to 0
$$
where the second map is given by $T_0, \ldots, T_n$. In this
case the statement above translates into the statement that
the elements
$$
T_{ij} = (0, \ldots, 0, T_j, 0, \ldots, 0, -T_i, 0, \ldots, 0)
\in M(1)_0
$$
generate the graded module $M(1)$ over $S$. We omit the details.
\end{proof}








